\section{Apéndice: formulario de revisión de cambios en plataformas comunitarias}

Este apéndice condensa los casos de esta clase en un proceso de revisión reutilizable, aplicable a cambios de alto impacto como políticas de contenido, API, funciones de pago, resúmenes generados por IA e interfaces para agent. La revisión no busca que todos los grupos lleguen a un acuerdo, sino exigir que el equipo haga públicos, antes del lanzamiento, las dependencias clave, los costos, los cambios de poder y las rutas de recuperación ante fallas. Cada punto debe tener un responsable, evidencia y condiciones de detención, para no tomar «la comunidad se adaptará» como supuesto por defecto.

\begin{table}[H]
\centering
\begin{tabular}{L{0.18\textwidth}L{0.33\textwidth}L{0.39\textwidth}}
\toprule
Dimensión de revisión & Pregunta obligatoria & Evidencia mínima \\
\midrule
Contrato con el usuario & Quién gana o pierde qué capacidades & Recorrido del usuario, diferencias de permisos, ejemplos de notificación \\
Poder en la comunidad & Cómo cambia el poder de moderadores, creadores, miembros y plataforma & Matriz de decisiones, rutas de apelación y de salida \\
Modelo económico & Quién asume los costos y quién obtiene los ingresos o el valor de los datos & Rangos de costo, nivel gratuito, reparto de ingresos y excepciones \\
Gobernanza de seguridad & Qué abusos y externalidades amplificará la nueva función & Modelo de amenazas, herramientas de moderación, proceso de escalamiento \\
Infraestructura & Qué acciones se conservan en picos de carga, fallas y reversiones & Pruebas de capacidad, orden de degradación, simulacros de recuperación \\
Ciclo de vida de los datos & Cómo se tratan los datos derivados tras la eliminación de contenido o la revocación de autorizaciones & Plazos de retención, limpieza de caché, registros de auditoría \\
\bottomrule
\end{tabular}
\caption{Todo cambio de alto impacto en una plataforma debe aprobar a la vez las revisiones de producto, gobernanza, economía, infraestructura y datos.}
\end{table}

\begin{table}[H]
\centering
\begin{tabular}{L{0.2\textwidth}L{0.32\textwidth}L{0.38\textwidth}}
\toprule
Etapa de lanzamiento & Indicadores adelantados & Condiciones de detención o reversión \\
\midrule
Pruebas internas & Tasa de permisos correctos, bloqueos erróneos, latencia y costo & Ampliación inexplicable de permisos o filtración de datos \\
Piloto en comunidades pequeñas & Adopción, quejas, horas de trabajo de los moderadores, tasa de apelaciones exitosas & Daño concentrado en grupos vulnerables o herramientas insuficientes \\
Ampliación del alcance & Tasa de migración del ecosistema, tasa de fallas, ingresos y oferta de contenido & Caída de la oferta principal, interrupción de clientes clave \\
Operación estable & Retención a largo plazo, coherencia de la gobernanza, recuperación de costos & Deuda en crecimiento constante sin un plan de corrección creíble \\
Recuperación de incidentes & Tiempo de recuperación, rezago acumulado, compensación y transparencia & Datos imposibles de conciliar o daños repetidos sin control \\
\bottomrule
\end{tabular}
\caption{Cada indicador debe ir acompañado de una condición de detención; de lo contrario, el monitoreo solo registrará las fallas sin cambiar las decisiones.}
\end{table}

\begin{warningbox}{Los tres grupos que con más facilidad se omiten en la revisión}
El primero son los desarrolladores que mantienen herramientas públicas sin tener un contrato comercial; el segundo, los moderadores voluntarios que asumen gran parte de la moderación sin figurar en la plantilla de personal; el tercero, los autores de contenido afectados por el entrenamiento con datos o por los resúmenes generados por IA, que ya no visitan directamente la plataforma. Si un cambio solo contabiliza a los clientes de pago y a los usuarios activos en las páginas, subestimará de forma sistemática las contribuciones y las pérdidas de estos grupos.
\end{warningbox}

\end{document}
